\documentclass[a4paper,12pt]{article}

\usepackage[a4paper,margin=16mm]{geometry}
\usepackage[german,main=french]{babel}
\usepackage{fontspec}
\usepackage{microtype}

\setmainfont{EBGaramond}
[
  Extension      = .otf,
  UprightFont    = *-Regular,
  ItalicFont     = *-Italic,
  BoldFont       = *-Bold,
  BoldItalicFont = *-BoldItalic,
  Numbers        = Lowercase
]

\usepackage{paracol}
\newcommand\chunks[2]{%
\begin{leftcolumn*}\begin{otherlanguage}{german}%
{#1}%
\end{otherlanguage}\end{leftcolumn*}%
\begin{rightcolumn}\begin{otherlanguage}{french}%
{#2}%
\end{otherlanguage}\end{rightcolumn}%
}

\newcommand\strophe[1]{\subsection*{\normalsize\normalfont\itshape #1}}

\begin{document}\thispagestyle{empty}
\setlength{\columnsep}{8mm}

\begin{paracol}{2}
  \chunks{\section*{O Haupt voll Blut und Wunden}}
         {\section*{Ô tête couverte de sang et de plaies}}

  \chunks{\noindent\raggedright\small Paul Gerhardt (1656), nach \textit{Salve caput cruentatum}}
         {\noindent\raggedright\small Paul Gerhardt (1656), d’après le \textit{Salve caput cruentatum}}

  \chunks{\noindent\raggedright\small Johann Sebastian Bach, \textit{Matthäus-Passion} (1727)}
         {\noindent\raggedright\small Jean-Sébastien Bach, \textit{Passion selon saint Matthieu} (1727)}

  \chunks{\strophe{Nr.~15 (Strophe 5)}
          \begin{verse}
          Erkenne mich, mein Hüter,\\
          Mein Hirte, nimm mich an!\\
          Von dir, Quell aller Güter,\\
          Ist mir viel Guts getan.\\
          Dein Mund hat mich gelabet\\
          Mit Milch und süßer Kost,\\
          Dein Geist hat mich begabet\\
          Mit mancher Himmelslust.
          \end{verse}}
         {\strophe{N\textsuperscript{o}~15 (strophe 5)}
          \begin{verse}
          Reconnais-moi, mon gardien,\\
          Mon berger, accueille-moi !\\
          De toi, source de tous les biens,\\
          J’ai reçu beaucoup de bien.\\
          Ta bouche m’a nourri\\
          De lait et de doux aliments,\\
          Ton esprit m’a comblé\\
          De mainte joie céleste.
          \end{verse}}

  \chunks{\strophe{Nr.~17 (Strophe 6)}
          \begin{verse}
          Ich will hier bei dir stehen;\\
          Verachte mich doch nicht!\\
          Von dir will ich nicht gehen,\\
          Wenn dir dein Herze bricht.\\
          Wenn dein Herz wird erblassen\\
          Im letzten Todesstoß,\\
          Alsdenn will ich dich fassen\\
          In meinen Arm und Schoß.
          \end{verse}}
         {\strophe{N\textsuperscript{o}~17 (strophe 6)}
          \begin{verse}
          Je veux me tenir ici près de toi ;\\
          Ne me méprise donc pas !\\
          Je ne veux pas te quitter\\
          Quand ton cœur se brise.\\
          Quand ton cœur pâlira\\
          Sous le dernier coup de la mort,\\
          Alors je veux te saisir\\
          Dans mes bras, sur mon sein.
          \end{verse}}

  \chunks{\strophe{Nr.~54 (Strophen 1 und 2)}
          \begin{verse}
          O Haupt voll Blut und Wunden,\\
          Voll Schmerz und voller Hohn,\\
          O Haupt, zu Spott gebunden\\
          Mit einer Dornenkron,\\
          O Haupt, sonst schön gezieret\\
          Mit höchster Ehr und Zier,\\
          Jetzt aber hoch schimpfieret,\\
          Gegrüßet seist du mir!
          \end{verse}}
         {\strophe{N\textsuperscript{o}~54 (strophes 1 et 2)}
          \begin{verse}
          Ô tête couverte de sang et de plaies,\\
          Pleine de douleur et d’outrages,\\
          Ô tête, par dérision ceinte\\
          D’une couronne d’épines,\\
          Ô tête, jadis bellement parée\\
          D’honneur et de gloire suprêmes,\\
          Mais à présent cruellement outragée,\\
          Je te salue !
          \end{verse}}

  \chunks{\begin{verse}
          Du edles Angesichte,\\
          Dafür sonst schrickt und scheut\\
          Das große Weltgewichte,\\
          Wie bist du so bespeit;\\
          Wie bist du so erbleichet!\\
          Wer hat dein Augenlicht,\\
          Dem sonst kein Licht nicht gleichet,\\
          So schändlich zugericht’?
          \end{verse}}
         {\begin{verse}
          Ô noble visage,\\
          Devant qui tremblent d’effroi\\
          Les grandes puissances du monde,\\
          Comme te voilà couvert de crachats ;\\
          Comme te voilà pâli !\\
          Qui donc a la lumière de tes yeux,\\
          Que nulle autre lumière n’égale,\\
          Si honteusement meurtrie ?
          \end{verse}}

  \chunks{\strophe{Nr.~62 (Strophe 9)}
          \begin{verse}
          Wenn ich einmal soll scheiden,\\
          So scheide nicht von mir,\\
          Wenn ich den Tod soll leiden,\\
          So tritt du denn herfür!\\
          Wenn mir am allerbängsten\\
          Wird um das Herze sein,\\
          So reiß mich aus den Ängsten\\
          Kraft deiner Angst und Pein!
          \end{verse}}
         {\strophe{N\textsuperscript{o}~62 (strophe 9)}
          \begin{verse}
          Quand un jour il me faudra partir,\\
          Alors ne te sépare pas de moi ;\\
          Quand il me faudra souffrir la mort,\\
          Alors avance-toi !\\
          Quand mon cœur sera\\
          Au comble de l’angoisse,\\
          Arrache-moi à mes angoisses\\
          Par ton angoisse et ta peine !
          \end{verse}}
\end{paracol}
\end{document}
